
This section reviews the two traditions the paper connects: probabilistic semantics, and
durable execution by logging and replay.

\subsection{Probabilistic Semantics}
\label{sec:bg-prob}

Probabilistic program semantics goes back at least to Kozen~\citep{kozen1981semantics}, who
interpreted programs as transformers of probability measures.
Borgstr\"{o}m et al.~\citep{borgstrom2016probabilistic} give distribution-based big-step and
small-step semantics for a probabilistic call-by-value $\lambda$-calculus and prove them
equivalent, together with a sampling-based semantics that agrees with them when integrated over
traces.
Our semantics is of the sampling kind: a run reads its random choices from a tape, and the
distribution over runs is induced by the distribution over tapes (Lemma~\ref{lem:tape}).
Fixing the tape is also what lets us compare runs in different modes directly, an instance of
the coupling method for relational reasoning about probabilistic
programs~\citep{barthe2015coupling}.
Probabilistic automata~\citep{segala1995probabilistic} add nondeterministic choice, resolved by
a scheduler, to probabilistic transitions; a single agent needs no scheduler, but concurrent
agents would (Section~\ref{sec:limitations}).

\subsection{Durable Execution}
\label{sec:bg-durable}

Durable-execution systems, Temporal~\citep{temporal2024docs} being the canonical production
example, give long-running workflows crash-recoverable semantics.
The mechanism is an \emph{event history}: an append-only log of all externally determined
values received by the workflow.
On restart, the workflow code is re-executed from the beginning, but any call that would
normally block on an external response instead returns the stored value from the log.
Workflow code must therefore be deterministic: re-run against the same history, it must make
the same decisions.
Calls to the outside world, including LLM calls, run as \emph{Activities} whose results are
recorded once and thereafter returned from the history, and other nondeterministic values such
as random numbers are obtained through replay-safe SDK operations that record
them~\citep{temporal2026sideeffects}.
This makes replay correct, but it gives no semantics to the \emph{distribution} of an
Activity's outcome, only to the single realized value.

The technique itself is classical.
Message-logging protocols for rollback recovery~\citep{strom1985optimistic,elnozahy2002survey}
assume that processes are \emph{piecewise deterministic}: execution is deterministic between
nondeterministic events, which are logged and replayed.
Our journal applies the same idea to an LLM agent, for which every model call is a
nondeterministic event, and our contribution is to state and prove its soundness inside a
probabilistic semantics.

\subsection{Comparison}

Table~\ref{tab:comparison} compares this work with the closest related systems;
Section~\ref{sec:related} discusses them in detail.

\begin{table}[ht]
\centering
\caption{Comparison with the closest related systems.
  \checkmark~= supported, $\times$~= not supported, $\sim$~= partial.}
\label{tab:comparison}
\fitwidth{%
\begin{tabular}{lccccc}
\toprule
\textbf{System} & \textbf{Multi-agent} & \textbf{Prob.\ semantics} & \textbf{Durable replay} & \textbf{Mechanized} & \textbf{Impl.} \\
\midrule
LLMbda~\citep{garby2026llmbda}
  & $\times$ & \checkmark & $\times$ & \checkmark{} (Lean) & \checkmark{} (verified interpreter) \\
MSC coord.~\citep{bollig2026msc}
  & \checkmark & $\times$ (black box) & $\sim$ (ZipperGen, informal) & \checkmark{} (Lean~4) & \checkmark{} (Python) \\
Temporal~\citep{temporal2024docs}
  & $\sim$ & $\times$ (det.\ required) & \checkmark & $\times$ & \checkmark \\
\textbf{This work}
  & $\times$ & \checkmark & \checkmark & \checkmark{} (Lean, pointwise) & \checkmark{} (Rust) \\
\bottomrule
\end{tabular}%
}
\end{table}
